%% kapa: exam=1 round=30 year=2019 period=1 form=A subject=civil_law subject_ko=민법 seq=15 number=15
%% source: Q-Net 감정평가사 기출문제 (cst00309) artlSeq=5205205 file=extracted/1008/5205205/2019년 제30회 감정평가사 1차시험 문제지/2019년 제30회 감정평가사 1차시험 문제지/2019년도 제30회 감정평가사 1차 1교시 A형.hwp
%% conversion: kordoc 4.24.51 -> tools/gen1.py (auto); review=proofread-sonnet; flags=none
\begin{question}{15}
甲으로부터 대리권을 수여받지 못한 乙은 甲의 대리인이라고 사칭하여 甲의 토지에 대해 丙과 매매계약을 체결하였다. 甲, 乙, 丙 사이의 법률관계에 관한 설명으로 옳은 것은? (다툼이 있으면 판례에 따름)
\begin{choices}
  \item 甲은 乙의 대리행위를 추인할 수 있으며, 그 추인은 乙이 아닌 丙에게 하여야 효력이 있다.
  \item 甲이 추인하지 않고 乙이 자신의 대리권을 증명하지 못한 경우, 乙은 자신의 선택에 좇아 선의ㆍ무과실인 丙에게 계약의 이행이나 손해배상 책임을 진다.
  \item 甲이 추인하면서 특별한 의사표시를 하지 않았다면 乙의 대리행위는 추인한 때로부터 甲에게 효력이 생긴다.
  \item 丙이 계약당시에 乙에게 대리권이 없다는 사실을 알았다면 철회권을 행사할 수 없다.
  \item 丙은 甲에게 상당한 기간을 정하여 추인 여부의 확답을 최고할 수 있으며, 甲이 그 기간 내에 확답을 발하지 아니하면 甲이 추인한 것으로 본다.
\end{choices}
\end{question}
\answer{4}
